\clearpage
\SourceLeaf{242}\pdfbookmark[1]{Dritter Abschnitt.}{liber-secundus-sectio-tertia}
\begin{center}
{\large DRITTER ABSCHNITT.}\par\medskip
\emph{Bestimmung einer Bahn, die beliebig vielen Beobachtungen möglichst genau genügt.}\par\bigskip
\hypertarget{native-242-article172}{}\pdfbookmark[2]{172.}{article-172}172.
\end{center}
Hätten die astronomischen Beobachtungen und die übrigen Zahlen, auf denen die Bahnberechnung beruht, absolute Präzision, so erhielte man auch die Bahnelemente, ob sie nun auf drei oder auf vier Beobachtungen beruhen, sofort mit absoluter Genauigkeit (sofern man freilich annimmt, dass die Bewegung genau nach den Gesetzen \textsc{Kepler}s erfolgt), und weitere Beobachtungen könnten sie daher nur bestätigen, niemals korrigieren. Da aber alle unsere Messungen und Beobachtungen nur Annäherungen an die Wahrheit sind und dasselbe für alle darauf beruhenden Rechnungen gelten muss, ist es das oberste Ziel aller Berechnungen konkreter Erscheinungen, der Wahrheit so nahe wie möglich zu kommen. Dazu müssen \emph{mehr} Beobachtungen, als zur Bestimmung der unbekannten Größen unbedingt erforderlich sind, in geeigneter Weise kombiniert werden. Diese Aufgabe lässt sich erst in Angriff nehmen, wenn die Bahn bereits näherungsweise bekannt ist. Sie ist dann so zu korrigieren, dass sie allen Beobachtungen \emph{möglichst genau} genügt. Mag dieser Ausdruck auch etwas unbestimmt erscheinen, so werden wir doch weiter unten die Grundsätze angeben, nach denen sich die Aufgabe streng und methodisch lösen lässt.

Nach höchster Präzision zu streben kann sich erst lohnen, wenn man der Bahnbestimmung gewissermaßen den letzten Schliff gibt. Solange hingegen Aussicht besteht, dass neue Beobachtungen bald zu weiteren Korrekturen \SourceLeaf{243}Anlass geben werden, empfiehlt es sich, je nach Lage der Dinge mehr oder weniger auf äußerste Präzision zu verzichten, sofern sich dadurch der Rechenaufwand erheblich verringern lässt. Wir werden versuchen, beiden Fällen gerecht zu werden.
